\subsection{Admissible topology of Poincaré groups}

Let B be a topological space and let \(a\) be a point of B. We say that a subgroup H of \(\pi_1(B, a)\) is admissible if every point \(b\) of B has a neighborhood V such that \(\gamma i_*(\delta)\gamma^{-1} \in \mathrm{H}\) for every \(\gamma \in \varpi_{a,b}(\mathrm{B})\) and every \(\delta \in \pi_1(V, b)\) , where \(i: V \to B\) is the canonical injection. If H is a normal subgroup of \(\pi_1(B, a)\) , it suffices, for H to be admissible, that this condition be verified for one path class \(\gamma \in \varpi_{a,b}(\mathrm{B})\) .

Proposition 4. — There exists a unique topology on \(\pi_1(B, a)\) compatible with its group structure, for which the admissible distinguished subgroups of \(\pi_1(B, a)\) form a fundamental system of neighborhoods of the identity element. For this topology, the open subgroups are exactly the admissible subgroups.

The following facts result from the definition of an admissible subgroup:

\begin{enumerate}
\def\labelenumi{\alph{enumi})}
\item
  The group \(\pi_1(B, a)\) is admissible;
\item
  A subgroup containing an admissible subgroup is admissible;
\item
  The intersection of a finite family of admissible subgroups is admissible.
\item
  For every admissible subgroup H of \(\pi_{1}(B,a)\) the intersection of the subgroups \(\gamma H\gamma^{-1}\) , when \(\gamma\) runs through \(\pi_{1}(B,a)\) , is admissible.
\end{enumerate}

In particular, the set of admissible normal subgroups is a filter base consisting of subgroups of \(\pi_1(B,a)\) satisfying the axiom (GV \(_{\text{III}}'\) ) of III, p. 4, whence the first part of the proposition according to TG, III, p. 5, example. The second part is then immediate (cf. TG, III, p. 7, corollary of prop. 4).

The topology on the group \(\pi_{1}(B,a)\) characterized in Proposition 4 is called the admissible topology.

Remarks. --- 1) If B \(_{0}\) denotes the arcwise connected component of \(a\) in B, the canonical isomorphism \(\pi_{1}(B_{0}, a) \to \pi_{1}(B, a)\) is an isomorphism of topological groups when these groups are equipped with the admissible topology.

\begin{enumerate}
\def\labelenumi{\arabic{enumi})}
\setcounter{enumi}{1}
\item
  Let B be a topological space. Let \(b\) and \(b'\) be points of B belonging to the same path-connected component, and let \(\gamma\) be an element of \(\varpi_{b,b'}(B)\) . It follows from the definition of an admissible subgroup that, for every admissible subgroup H of \(\pi_1(B,b')\) , the subgroup \(\gamma H\gamma^{-1}\) of \(\pi_1(B,b)\) is admissible. Consequently, the isomorphism \(u_\gamma \colon \delta \mapsto \gamma \delta \gamma^{-1}\) of \(\pi_1(B,b')\) onto \(\pi_1(B,b)\) (cf. III, p. 292) is a homeomorphism when these groups are equipped with the admissible topologies.
\item
  Let A and B be topological spaces, let \(a\) be a point of A and let \(f\colon\mathrm{A}\to\mathrm{B}\) be a continuous map. Set \(b=f(a)\) . If H is an admissible subgroup of \(\pi_{1}(\mathrm{B},b)\) , its inverse image under the homomorphism \(\pi_{1}(f,a)\) is an admissible subgroup of \(\pi_{1}(\mathrm{A},a)\) . Consequently, the group homomorphism \(\pi_{1}(f,a)\colon\pi_{1}(\mathrm{A},a)\to\pi_{1}(\mathrm{B},b)\) is continuous, when these groups are equipped with the admissible topology.
\end{enumerate}

PROPOSITION 5. --- Let B be a locally path-connected topological space and let a be a point of B. For a subgroup H of \(\pi_1(\mathrm{B},a)\) to be admissible, it is necessary and sufficient that there exist a covering (E,p) of B and a point \(x\in \mathrm{E}_a\) such that \(\mathrm{H} = p_{*}(\pi_{1}(\mathrm{E},x))\) .

Let \((\mathrm{E},p)\) be a covering of B and let x be a point of \(E_{a}\) ; set \(\mathrm{H}=p_{*}(\pi_{1}(\mathrm{E},x))\) and let us show that it is an admissible subgroup of \(\pi_{1}(\mathrm{B},a)\) . Let b be a point of B and let V be a neighborhood of b such that \(\mathrm{E}_{\mathrm{V}}=(p^{-1}(\mathrm{V}),p_{\mathrm{V}})\) is a trivializable covering of V. Let \(\gamma\in\varpi_{a,b}(\mathrm{B})\) ; we shall prove that for every element \(\delta\in\pi_{1}(\mathrm{V},b)\) , the path class \(\gamma\delta\gamma^{-1}\) belongs to H.

By III, p. 301, prop. 3, there exists a unique strict homotopy class \(\gamma'\) of a path with origin \(x\) in E such that \(p_*(\gamma') = \gamma\) . Let \(y\) be the terminal point of \(\gamma'\) ; we have \(p(y) = b\) (III, p. 302, prop. 4). Let then \(\delta'\) be the unique path class with origin \(y\) in E such that \(p_*(\delta') = \delta\) . Since \(\mathrm{E_V}\) is trivializable, \(\delta'\) is the class of a loop at \(y\) . Then, \(\gamma'\delta'(\gamma')^{-1}\) is the class of a loop at \(a\) in E whose image under \(p_*\) is the class \(\gamma\delta\gamma^{-1}\) , which is what had to be proved.

Conversely, let H be an admissible subgroup of \(\pi_1(B, a)\) . Let \(\lambda_a(B)\) be the quotient of the space \(\Lambda_a(B)\) of paths with origin \(a\) for the strict homotopy relation; since two strictly homotopic paths have the same endpoint, the endpoint map \(e: \Lambda_a(B) \to B\) defines by passage to the quotient a map \(\varepsilon: \lambda_a(B) \to B\) . The composition of path classes endows the set \(\lambda_a(B)\) with a left action of the group \(\pi_1(B, a)\) . Equip it with the action of the group H obtained by restriction, and denote by \(H\backslash\lambda_a(B)\) the set of its orbits.

The map \(\varepsilon\colon\lambda_{a}(B)\to B\) induces, by passing to the quotient, a map \(q\colon H\backslash\lambda_{a}(B)\to B\) . The composition of path classes endows the set \(H\backslash\lambda_a(B)\) with a right action of the groupoid \(\varpi(B)\) with respect to the map \(q\) .

This operation has no local monodromy. Indeed, let \(b\) be a point of B, and let V be a neighborhood of \(b\) such that \(\gamma i_{*}(\pi_{1}(V,b))\gamma^{-1}\subset H\) for every path class \(\gamma\) connecting \(a\) to \(b\) in B, where \(i\) denotes the inclusion of V into B. Since B is locally path-connected, one may further assume that V is path-connected. Let \(c\in V\) , let \(\delta\) be the class in \(\pi_1(B,c)\) of a loop at \(c\) contained in V, and let \(\delta'\) be an element of \(\lambda_a(B)\) such that \(\varepsilon (\delta ') = c\) . Let \(\delta''\) be an element of \(\varpi_{c,b}(V)\) and let us set \(\gamma = \delta '\delta ''\) . By definition of V, the element \(\delta '\delta (\delta ')^{-1} = \gamma ((\delta '')^{-1}\delta \delta '')\gamma^{-1}\) of \(\pi_1(B,a)\) belongs to H. Then, \(\mathrm{H}\delta^{\prime}\cdot \delta = \mathrm{H}\delta^{\prime}\) ; this proves that \(\pi_1(V,c)\) acts trivially on the set \(q^{-1}(c)\) .

By III, p. 313, prop. 3, there exists a unique topology on \(\mathrm{H}\backslash\lambda_{a}(\mathrm{B})\) for which \(q\) is continuous and the B-space \((\mathrm{H}\backslash\lambda_{a}(\mathrm{B}),q)\) is a covering such that the canonical action of \(\varpi(\mathrm{B})\) on this covering is the action defined above. By assertion b) of Theorem 1 of III, p. 305, the group \(q_{*}(\pi_{1}(\mathrm{H}\backslash\lambda_{a}(\mathrm{B}),\mathrm{H}))\) is equal to H.

We will say that an operation of the group \(\pi_1(B, b)\) on a set X is admissible if the kernel of the canonical map \(\pi_1(B, b) \to \mathfrak{S}_{\mathrm{X}}\) is an open subgroup of \(\pi_1(B, b)\) . This amounts to saying that the homomorphism \(\pi_1(B, b) \to \mathfrak{S}_{\mathrm{X}}\) is continuous if the group \(\mathfrak{S}_{\mathrm{X}}\) is equipped with the discrete topology. The map \(\pi_1(B, b) \times X \to X\) is then continuous, when X is endowed with the discrete topology. Conversely, consider a continuous action of \(\pi_{1}(B,b)\) on a discrete space X. Let x be a point of X; its stabilizer is an open subgroup H of \(\pi_{1}(B,b)\) by hypothesis. For \(\gamma\in\pi_{1}(B,b)\) the fixator of \(\gamma\cdot x\) is the subgroup \(\gamma H\gamma^{-1}\) , so that the subgroup of \(\pi_{1}(B,b)\) fixing each element of the orbit of x is the intersection of the subgroups \(\gamma H\gamma^{-1}\) , for \(\gamma\) running through \(\pi_{1}(B,b)\) . It is an open subgroup because the distinguished open subgroups of \(\pi_{1}(B,b)\) form a basis for its topology (III, p. 315, prop. 4). This implies that a continuous action of \(\pi_{1}(B,b)\) on a discrete space X is admissible if it is transitive or, more generally, if the set of orbits of the elements of X is finite.

For any discrete group \(G\), any principal \(G\)-covering \(E\) over \(B\), and any point \(x \in E_{b}\) , recall that the map \(h_{(\mathrm{E},x)}\) from \(\pi_{1}(\mathrm{B},b)\) to G which, to \(\gamma \in \pi_{1}(\mathrm{B},b)\) , associates the unique element \(g \in G\) such that \(x \cdot g = x \cdot \gamma^{-1}\) is a group homomorphism (III, p. 306, prop. 5).

PROPOSITION 6. --- Let B be a topological space and let b be a point of B. Endow the group \(\pi_1(B, b)\) with the admissible topology.

\begin{enumerate}
\def\labelenumi{\alph{enumi})}
\item
  For any discrete group G, any principal covering E of B with group G, and any \(x \in E_{b}\) , the homomorphism \(h_{(E,x)}\) is a continuous homomorphism of topological groups.
\item
  For every covering E of B, the canonical action of \(\pi_{1}(B,b)\) on the fiber \(E_{b}\) is admissible.
\end{enumerate}

It suffices to prove the second assertion. Let K be the set of elements \(k \in \pi_{1}(B, b)\) such that \(x \cdot k = x\) for all \(x \in E_{b}\) . Let us prove that K is an admissible subgroup of \(\pi_{1}(B, b)\) .

Let \(b'\) be a point of B and let \(V'\) be a neighborhood of \(b'\) in B above which the covering E is trivializable. Let \(\gamma\) be the class of a path c connecting \(b\) to \(b'\) in B and let \(c'\) be the unique path starting at x in E that lifts c. For all \(\delta \in \pi_1(V', b')\) and every \(x' \in E_{b'}\) , we have \(x' \cdot \delta = x'\) . Consequently, for \(x \in E_b\) and \(\delta \in \pi_1(V', b')\) , we have

\[
x \cdot \gamma \delta \gamma^ {- 1} = ((x \cdot \gamma) \cdot \delta) \cdot \gamma^ {- 1} = (x \cdot \gamma) \cdot \gamma^ {- 1} = x,
\]

so that \(\gamma\delta\gamma^{-1}\) belongs to K. In other words, K is an admissible subgroup, whence the proposition.

PROPOSITION 7. --- Let B be a connected and locally path-connected topological space, and let b be a point of B. Endow the group \(\pi_1(B, b)\) with the admissible topology.

\begin{enumerate}
\def\labelenumi{\alph{enumi})}
\item
  Let G be a discrete topological group. For every continuous homomorphism \(f\colon \pi_{1}(\mathrm{B},b)\to \mathrm{G}\) , there exists a principal covering E of B with group G, and a point \(x\) of \(\mathrm{E}_{b}\) such that \(h_{(\mathrm{E},x)} = f\) .
\item
  For every discrete topological space F equipped with an admissible right action of the group \(\pi_{1}(B,b)\), there exists a covering E of B such that the \(\pi_{1}(B,b)\)-sets F and \(E_{b}\) are isomorphic.
\end{enumerate}

Let us prove \(a\) ). Let \(f\) be a continuous homomorphism of \(\pi_1(B, b)\) into a discrete group G. Its kernel is an open normal subgroup K of \(\pi_1(B, b)\) ; denote by H the group \(\pi_1(B, b)/K\) and \(\overline{f} \colon H \to G\) the group homomorphism induced by \(f\) by passing to the quotient. By III, p. 316, prop. 5, there exists a connected covering E' of B such that the subgroup K is the stabilizer of every point of the fibre \(E'_{b}\) ; the covering E' is principal with group H = \(\pi_1(B, b)/K\) . Let \(x'\) be a point of \(E'_{b}\) ; the homomorphism \(h_{(E', x')} \colon \pi_1(B, b) \to H\) is surjective with kernel K (III, p. 306, prop. 5). Let then \(\varphi \colon H \to G\) be the unique homomorphism of groups such that \(\varphi \circ h_{(E', x')} = f\) . The associated covering \(E = E' \times^{H} G\) of B is principal with group G, and one has \(h_{(E, x)} = \varphi \circ h_{(E', x')} = f\) (III, p. 307, example 2). This proves assertion \(a\) ).

Let us prove \(b\) ). Let F be a set equipped with an admissible right action of the group \(\pi_1(B,b)\) . Denote by \(f: \pi_1(B,b) \to \mathfrak{S}_{\mathrm{F}}\) this action and endow the group \(\mathfrak{S}_{\mathrm{F}}\) with the discrete topology. By \(a\) ), there exists a principal covering E of B with group \(\mathfrak{S}_{\mathrm{F}}\) , and a point \(x \in \mathrm{E}\) such that the homomorphism \(h_{(\mathrm{E},x)}\) is equal to \(f\) . The canonical action of the group \(\pi_1(B,b)\) on the fibre over \(b\) of the associated covering \(\mathrm{E} \times^{\mathfrak{S}_{\mathrm{F}}} \mathrm{F}\) of B is identified with the action of \(\pi_1(B,b)\) on F (III, p. 306, example 1).

Remarks. --- 4) The admissible topology on \(\pi_1(B, b)\) is the coarsest topology for which the action of \(\pi_1(B, b)\) on \(E_b\) is continuous for every connected covering E of B (prop. 6 of III, p. 318 and prop. 7 of III, p. 318). If E is a connected covering of B and \(x\) a point of the fibre \(E_b\) , the map \(c' \mapsto c'(1)\) of \(\Lambda_x(E)\) in E is continuous. Consequently, the map \(c \mapsto x \cdot c\) of \(\Lambda_b(B)\) in E is continuous (III, p. 302, cor. 2 of prop. 3). By restriction to \(\Omega_b(B)\) and passage to the quotient, one deduces that the map \(\gamma \mapsto x \cdot \gamma\) is continuous when one endows the group \(\pi_1(B, b)\) with the quotient topology of \(\Omega_b(B)\) . The admissible topology on \(\pi_{1}(B,b)\) is therefore coarser than the quotient topology of the topology of compact convergence. It can be strictly coarser (III, p. 337, exerc. 7).

\begin{enumerate}
\def\labelenumi{\arabic{enumi})}
\setcounter{enumi}{4}
\tightlist
\item
  Let B be a connected and locally path-connected topological space, let \(a\) be a point of B and let H be a subgroup of \(\pi_1(B,a)\) . By the preceding remark, if H is admissible, it is also open for the quotient topology of the topology of compact convergence. Conversely, suppose that H is a distinguished subgroup of \(\pi_1(B,a)\) which is open for the quotient topology of the topology of compact convergence. Let \(x\) be a point of B and let \(\gamma \in \varpi_{a,x}(B)\) . The subgroup \(\gamma^{-1}H\gamma\) of \(\pi_1(B,x)\) is still open for the quotient topology of the topology of compact convergence (III, p. 293, remark 3). There therefore exists a neighborhood V of \(x\) in B such that \(\gamma^{-1}H\gamma\) contains the class of every loop at \(x\) whose image is contained in V. We thus have \(\gamma \pi_1(V,x)\gamma^{-1} \subset H\) ; since H is distinguished, this implies that H is an admissible subgroup of \(\pi_1(B,a)\) .
\end{enumerate}

